\documentclass{article}
%\usepackage[a4paper, total={12cm, 12cm}]{geometry}
\usepackage{graphicx} % Required for inserting images
\usepackage[utf8]{inputenc}
\usepackage{amsmath}
\usepackage{url}

\usepackage{biblatex}
\usepackage{adjustbox}
\addbibresource{forrasok.bib} %Import the bibliography file

% =============================================================
% User Specific Macros
% --------------------

%% Karnaught-tábla
\newcommand{\igaz}{\textit{t}}
\newcommand{\hamis}{\textit{h}}

%% Logikai muveletek
\newcommand{\cland}[2]{{#1} \land {#2}%
}
\newcommand{\clor}[2]{{#1} \lor {#2}%
}
\newcommand{\csup}[2]{{#1} \supset {#2}%
}
\newcommand{\csub}[2]{{#1} \subset {#2}%
}
% log. eqivalencia
\newcommand{\cleqv}[2]{{#1} \leftrightarrow {#2}%
}
\newcommand{\cneg}[2]{{#1} \neg {#2}%
}

%
\newcommand{\AND}{\mathrm{AND}}
\newcommand{\OR}{\mathrm{OR}}
\newcommand{\NAND}{\mathrm{NAND}}
\newcommand{\XOR}{\mathrm{XOR}}
\newcommand{\NOT}{\mathrm{NOT}}

% =============================================================

\title{Aarhus University \\ Cryptographic Computing \\ Theoretical Assignments }
\author{ Fanni Zsanett Ferenczi, \\ Miklos Daniel Vadasz}
\date{August 2025}

\begin{document}

\maketitle

%\include{info.tex}



\section{Theoretical Assignment}

\subsection{Card-based secure two party protocol}

4 ordinary cards with identical backs, fronts are either red (R) or blue (B).

Each player encodes their bit using two cards placed face-down in a row (left→right) private to them:

\begin{itemize}
    \item To encode bit 0, place the pair RB (from left top right).
    \item To encode bit 1, place the pair BR (from left top right).
\end{itemize}
    
So each player uses two cards; after both have placed, there are 4 face-down cards in a row: Alice’s two, then Bob’s two.


Protocol steps:

    \begin{itemize}
        \item Alice privately places two face-down cards on the table in order according to her bit (0 → RB, 1 → BR).
        \item Bob privately places his two face-down cards immediately to the right of Alice’s (again 0 → RB, 1 → BR)
        \item Now the row is four face-down cards: positions 1..4 (Alice1, Alice2, Bob1, Bob2).
        \item Reveal the two middle cards (positions 2 and 3 of the row) face up (leave positions 1 and 4 face-down).
    \end{itemize}
    
If the two revealed middle cards are RB from left to right, output 0. Otherwise (any other revealed pair: BR or RR or BB) output 1. 

\begin{table}[h!]
    \centering
    \begin{tabular}{|c|c||c|c|c|c||c|}
         \hline\hline
       \textit{Alice}  & \textit{Bob} & \textbf{A1} & \textbf{A2} & \textbf{B1} & \textbf{B2} & Output \\
         \hline 
       1  &1 & R & B & R &B &1 \\
         \hline 
       1  & 0 &R & B  & B& R &1 \\
         \hline 
       0  & 1 &B & R& R &B &1 \\
         \hline 
       0   & 0 &B &R &B &R &0 \\
        \hline
    \end{tabular}
    \caption{4-card protocol}
    \label{tab:4cardproto}
\end{table}

\subsection{Blood type compatibility}

%Submit description of the function $f$ as a truth table and as a Boolean formula.

We will use the following secure two-party computation scenario as a running example for all the assignments: 

\textit{Alice} has blood type $x$, Bob has blood type $y$. We want Alice to learn $z=f(x,y)$, where $f$ is the function that outputs $1$ if a person with blood type $x$ can receive a blood transfusion from a person with blood type $y$ (and $0$ otherwise). Bob should learn nothing.

%% TTP
\subsubsection{Description of the function $f$ as a truth table}

The compatibility matrix (truth table)\cite{au.mpc.a1.bloodtype} between the $8$ blood types can be represented in \textit{Table \ref{tab:bloodtypetruth}} as follows:

% Table head: donor, column := recipient
% O- := Universal Donor
% AB+ := Universal recipient 

% x := recipient
% y := donor


\begin{table}[h!]
    \centering
    \begin{tabular}{|l|c|c|c|c|c|c|c|c|}
    \hline
    \textit{Donor:} & $0-$ & $0+$ & $A-$ & $A+$ & $ B-$ & $B+$ & $AB-$ & $AB+$  \\
      \hline\hline
    \textit{Receiver:} $0-$  & 1 & 0 & 0 &  0 & 0 & 0 & 0  &  0 \\
      \hline
    \textit{Receiver:}  $0+$  & 1 & 1 & 0 & 0 &  0 & 0 & 0   & 0  \\
      \hline
    \textit{Receiver:}  $A-$ &1  & 0 & 1 &  0 & 0 & 0  &  0 &  0 \\
      \hline
    \textit{Receiver:}  $A+$ & 1 & 1 & 1 & 1 & 0 &  0 & 0  & 0  \\
      \hline
     \textit{Receiver:}  $ B-$ &  1& 0 & 0 & 0 &  1 &  0 & 0 &  0 \\
      \hline
      \textit{Receiver:} $B+$ & 1 & 1 & 0 & 0 & 1 &1 & 0  & 0  \\
      \hline
     \textit{Receiver:}  $AB-$ & 1  & 0 &1 & 0 &1 & 0 & 1 &  0 \\
      \hline
     \textit{Receiver:} $AB+$  &1 & 1 &  1 & 1 & 1 & 1 & 1 &  1 \\
      \hline
    \end{tabular}
    \caption{Blood type compatibility Matrix (truth table)}
    \label{tab:bloodtypetruth}
\end{table}

\subsection{Description of $f$ as Boolean Formula}

In the second case, we can also represent $f$ with the aid of boolean operators (AND, OR, XOR, NOT, etc.). 

For easier readability let $F_i$  represents $i$-th row vector of \textit{Table \ref{tab:bloodtypetruth}}.

Thus it yields us with the following representation: 

\begin{table}[h!]
    \centering
$
\begin{array}{c||l}
\hline
F_1 & O{-} \\
\hline
F_2 & O{-} \lor~~\  O{+} \\
\hline
F_3 & O{-} \lor~~\  A{-} \\
\hline
F_4 & O{-} \lor~~\  O{+} \lor~~\  A{-} \lor~~\  A{+} \\
\hline
F_5 & O{-} \lor~~\  B{-} \\
\hline
F_6 & O{-} \lor~~\  O{+} \lor~~\  B{-} \lor~~\  B{+} \\
\hline
F_7 & O{-} \lor~~\  A{-} \lor~~\  B{-} \lor~~\  AB{-} \\
\hline
F_8 & O{-} \lor~~\  O{+} \lor~~\  A{-} \lor~~\  A{+} \lor~~\  B{-} \lor~~\  B{+} \lor~~\  AB{-} \lor~~\  AB{+} \\
\hline
\end{array}
$    
    \caption{Rewriting the blood type \textit{Truth Table} with boolean operators}
    \label{tab:secondpoz}
\end{table}


Furthermore we can simplify our notation using $R_{\_}$ denoting the \textit{Receiver} and $D_{\_}$ as \textit{Donor}. The underscore ($\_$) denotes the placeholder for given blood type based on notation used within \textit{Table \ref{tab:bloodtypetruth}, \ref{tab:secondpoz}}.


%% 1st assignment 3 bit represantation alaphoz>
\[
\begin{aligned}
f =\ & (R_{O-} \land~~\ D_{O-}) \\
\lor~~\ & (R_{O+} \land~~\ (D_{O-} \lor~~\ D_{O+})) \\
\lor~~\ & (R_{A-} \land~~\ (D_{O-} \lor~~\ D_{A-})) \\
\lor~~\ & (R_{A+} \land~~\ (D_{O-} \lor~~\ D_{O+} \lor~~\ D_{A-} \lor~~\ D_{A+})) \\
\lor~~\ & (R_{B-} \land~~\ (D_{O-} \lor~~\ D_{B-})) \\
\lor~~\ & (R_{B+} \land~~\ (D_{O-} \lor~~\ D_{O+} \lor~~\ D_{B-} \lor~~\ D_{B+})) \\
\lor~~\ & (R_{AB-} \land~~\ (D_{O-} \lor~~\ D_{A-} \lor~~\ D_{B-} \lor~~\ D_{AB-})) \\
\lor~~\ & (R_{AB+} \land~~\ (D_{O-} \lor~~\ D_{O+} \lor~~\ D_{A-} \lor~~\ D_{A+} \lor~~\ D_{B-} \lor~~\ D_{B+} \lor~~\ D_{AB-} \lor~~\ D_{AB+}))
\end{aligned}
\]

\subsubsection*{ Compatibility matrix using 3-bit representation }

In case we are using one bit to represent modifier $+,-$ (1 if it represents $+$, and $-$ if it represents $0$)  we can further reformulate our \textit{Table \ref{tab:bloodtypetruth}}, leading us with a more easier approach to simply our boolean function representation of $f$. 

With encoding (3 bits per type) given the expression:  for any blood type $t$ let its 3-bit encoding be $t \rightarrow{ (a,b,r)}$ where: 
\begin{itemize}
    \item  $a = 1$ if the blood has the A antigen\cite{au.mpc.a1.bloodtype.theory} (types A or AB), otherwise $a = 0$,
    \item $b = 1$ if the blood has the B antigen\cite{au.mpc.a1.bloodtype.theory} (types B or AB), otherwise $b = 0$,
    \item $r = 1$ if Rh-positive ($+$), otherwise $r = 0$ for Rh-negative ( $-$ ).
\end{itemize}

\begin{center}
\begin{table}[!h]
\centering
\begin{tabular}{|c||c|}
\hline
\textbf{Blood Type} & \textbf{Binary representation: $A|B|\mathrm{\{+,-\}}$ } \\
\hline\hline 
O$-$ & 000  \\
O$+$ & 001  \\
A$-$ & 100 \\
A$+$ & 101 \\
B$-$ & 010 \\
B$+$ & 011  \\
AB$-$ & 110  \\
AB$+$ & 111 \\
\hline
\end{tabular}
\caption{Blood Type representation with 3-bit encoding scheme}
\label{tab:3bitreprz}
\end{table}
\end{center}

Under this encoding the eight types represented as the following \textit{Table \ref{tab:3bitreprz}}: 

\subsubsection*{Simplified boolean formula $f$}
\begin{itemize}
    \item Let recipient bits be $(x_A, x_B, x_R)$ and donor bits be $(y_A, y_B, y_R)$
    \item  A recipient of type $x$ can receive from a donor of type $y$ exactly when the donor does not have any antigen that the recipient lacks (for A and B) and the donor is not $Rh+$ while the recipient is $Rh-$
\end{itemize}
    
Thus, written with boolean operators we can express $f$ as follows:

\begin{equation}
    f(x,y) = (x_A \lor \neg y_A) \land (x_B \lor \neg y_B) \land (x_R \lor \neg y_R)
\end{equation}

\subsubsection*{Quick checks}

Finally we are able to check our formalization. 

\begin{itemize}
    \item  Donor O- (0,0,0) to any recipient: every $\neg y_{*}$ is true, so $f = 1$ 
    \item  Recipient AB+ (1,1,1) from any donor: every $x_{*}$  is 1, so $f = 1$
    \item  Recipient A+ (1, 0, 1) from donor O+ (0, 0, 1): $f = 1$, as expected 
    \item  Recipient A+ (1,0,1) from donor B+ (0,1,1): the second clause is 0, so $f = 0$ (incompatible), as expected
\end{itemize}


% Programming assignments .tex-be a leiras + formalizmus => git repo
% Project proposal + final project  =>  Masik overleaf project

%\include{prog.tex}




\printbibliography
\end{document}


